%\author{Javier Alejandro Acevedo Barroso}
%\date{July 2013}
%\title{ Thesis submitted for the degree of Physicist}
%\publishers{Universidad de los Andes}

\begin{titlepage}
   \begin{center}
      \Large\textbf{SIMULATING COLLISIONAL DARK}\\
      \vspace{-4mm}\Large\textbf{MATTER USING A LATTICE}\\
      \vspace{-4mm}\Large\textbf{BOLTZMANN METHOD}\\
\vspace{15mm}
      \Large\textbf{Javier Alejandro}\\
      \vspace{-5mm}\Large\textbf{Acevedo Barroso}\\
\vspace{11mm}      
      \Large\textbf{Advisor:}\\
      \vspace{-5mm}\Large\textbf{Dr. Jaime E. Forero-Romero}\\
\vspace{15mm}      
     \Large\text{A monograph presented for the degree of Physicist.}\\
\vspace{8mm}           
      {\includegraphics[scale=0.11]{imag/logoUniandes.png}}\\
\vspace{3mm}           
      \Large\text{Departamento de Física}\\
      \vspace{-4mm}\Large\text{Facultad de Ciencias}\\
      \vspace{-4mm}\Large\text{Universidad de los Andes}\\
      \vspace{-4mm}\Large\text{Bogotá, Colombia}\\
      \vspace{-1mm}\Large\text{November, 2018}
   \end{center}
\end{titlepage}

\thispagestyle{empty}
\begin{dedication}
\Huge\textit{Dedicado a}\\
%\vspace{-1mm}\large\text{Satán}\\
\vspace{-1mm}\large\text{Mi mamá, pues sin su paciencia y consejo nunca lo habría logrado.}\\
\vspace{-1mm}\large\text{Mi papá, quién me enseñó el valor del arte, la ciencia y la amistad.}\\
\vspace{-1mm}\large\text{Susana, por la inolvidable vida que hemos compartido.}\\
\end{dedication}


\newpage
\thispagestyle{empty}
\begin{center}
      \Large\textbf{Simulating Collisional Dark Matter}\\
      \vspace{-4mm}\Large\textbf{Using a Lattice Boltzmann Method}\\
      %\vspace{-4mm}\Large\textbf{BOLTZMANN METHOD}\\
\vspace{07mm}
      \Large\textbf{Javier Alejandro Acevedo Barroso}\\
      \vspace{-4mm}\Large{November, 2018}\\
      
\vspace{07mm}
      \Large\textbf{Abstract}\\
      \vspace{05mm}\small\justify{
Usually, dark matter is simulated with N-body schemes that sample the phase space in order to solve the Poisson-Vlasov equation. %\cite{2012PDU150K}.
These kind of simulations have been essential for the development of modern cosmology and the characterization of dark matter halos. With the development of particle physics, we ultimately expect dark matter to be a particle outside of the standard model of physics. Recent measurements on the aftermath of galaxy cluster collisions allow us to constrain the value of the thermally averaged cross section $\crosssection$, thus motivating the development of dark matter \emph{collisional} simulations.\\
On the other hand, Lattice-Boltzmann simulations have been widely used to recreate increasingly complex fluids and boundary conditions, nonetheless, the usual Lattice-Boltzmann scheme does not simulate the entirety of the velocity space, but simply a small number of adventive velocities. \\
Inspired by the work of Philip Mocz, Sauro Succi \cite{integerLatticeDynamics}, and Sebastian Franco \cite{franco}, in which a Lattice-Boltzmann scheme is used to simulate the phase space of a \emph{collisionless} one dimensional dark matter fluid. We implement a Lattice-Boltzmann simulations of the phase space of a \emph{collisional} one, two, and three dimensional dark matter fluid. For the collisional step, we use the BGK approximation\cite{1954PhRv...94..511B} \cite{asinari} modeled by a relaxation time $\tau$ chosen accordingly to recent measurements of the $\crosssection$.}\\
%      \vspace{-4mm}\Large\textbf{Using a Lattice Boltzmann Method}\\	
\end{center}

\newpage
\thispagestyle{empty}
\begin{center}
      \Large\textbf{Simulando Materia Oscura Colisional}\\
      \vspace{-4mm}\Large\textbf{Utilizando un Método de Lattice Boltzmann}\\
      %\vspace{-4mm}\Large\textbf{BOLTZMANN METHOD}\\
\vspace{07mm}
      \Large\textbf{Javier Alejandro Acevedo Barroso}\\
      \vspace{-4mm}\Large{Noviembre, 2018}\\
      
\vspace{07mm}
      \Large\textbf{Resumen}\\
      \vspace{05mm}\small\justify{
Usualmente la materia oscura se simula con métodos de N cuerpos, en los que se se muestrea el espacio de fase con el fin de resolver la ecuación de Poisson-Vlasov.\\
Este tipo de simulaciones ha sido esencial tanto para el desarrollo de la cosmología moderna, como para la caracterización de los halos de materia oscura. Con el desarrollo de la física de particulas, esperamos que en últimas la materia oscura sea una particula fuera del modelo estandar. Mediciones recientes del resultado de colisiones entre cúmulos galácticos, permiten dar límites para los valores del promedio térmico de la sección transversal $\crosssection$, motivando el desarrollo de simulaciones colisionales de materia oscura.\\
Por otro lado, las simulaciones de Lattice-Boltzmann han sido ampliamente utilizadas para recrear fluidos cada vez más complejos. Sin embargo, la implementación usual no simula el espacio de velocidades en sí, sino un pequeño número de velocidades adventivas.\\
Inspirados en el trabajo realizado por Philip Mocz, Sauro Succi \cite{integerLatticeDynamics}, y Sebastián Franco \cite{franco}, en donde se simula el espacio de fase de un fluido \emph{no colisional} de materia oscura en una dimensión. Implementamos un método de Lattice-Boltzmann para simular un fluido \emph{colisional} de materia oscura en una, dos y tres dimensiones. Para el paso colisional usamos la aproximación BGK\cite{1954PhRv...94..511B} \cite{asinari} modelada por un tiempo de relajación $\tau$, seleccionado de acuerdo a mediciones recientes de $\crosssection$.}\\
%      \vspace{-4mm}\Large\textbf{Using a Lattice Boltzmann Method}\\	
\end{center}
\thispagestyle{empty}
%      \vspace*{15mm}\Huge\textbf{Agradecimientos}\\
\chapter*{Agradecimientos}
      \normalsize{Agradezco profundamente a Jaime Forero, por su guía, paciencia y visión en el desarrollo de esta monografía. A Sebastián Franco, por sentar muy bien las bases de este proyecto. Y a todos los se convirtieron en mis amigos en estos años de estudio. }\\
